\subsection{Diagram 5}\label{sec:d5}

Diagram~5 is the isolated $zw$-triangle $345$ alone: its only strokes are the six strokes of
the three triangle edges, and no vertex is circled \cite[\S7.4 and Fig.~11]{AK}.
We write $T=\{3,4,5\}$ and call $1$ and $2$ the \emph{outer bodies}.  For this diagram
\cite[\S\S8.4--8.5]{AK} record only relation (23), and \cite[\S9]{AK} avoid estimating the
distances that carry no edge.
We build no chart.  We classify the orders of the outer bodies relative to the triangle and
decide by exact linear programming which are compatible with the equations on a curve through
a real point, where $U=I$ has a nonzero constant value.

\begin{theorem}\label{thm:d5}
Let $m\in\Rp^5$.  No Diagram 5 germ with an unbounded four-product lies on a curve through a
real point: if an irreducible curve $\Gamma\subset A_m$ carries such a germ, then $\Gamma$
contains no real point.
\end{theorem}

This is part~(b) of Theorem~\ref{thm:exclusion} for $r=5$.  Germs with all four-products bounded
are not excluded, and need not be (Proposition~\ref{prop:reduction}); they do occur on curves
through real points, e.g.\ at the collision end of the modified Roberts continuum at the masses
$(4,4,4,4,1)$ \cite[p.~543]{AK}, which is real but not positive.

\subsubsection{The frame}\label{d5-frame}
Let $\gamma$ be a Diagram 5 germ, numbered as in \cite[Fig.~11]{AK} and renormalized as in
Lemma~\ref{lem:germ}, with orders measured in units of $\ord\epsilon$, so that $\mathcal Z$ and
$\mathcal W$ have least order $-2$.  After a constant reparametrization the Puiseux parameter
$t$ has $\epsilon\approx|t|$, and we write orders in $t$.
By Fact~\ref{fact:rules}: $-2\le\ord r_{kl}\le1$, with $\ord r_{kl}=1$ exactly for $zw$-edges, and
$-2\le\ord z_{kl},\ord w_{kl}\le2$ (Estimate~1); $\ord z_{kl}=\ord w_{kl}=1$ and
$\ord\delta_{kl}=-1$ on $T$ (Estimate~2); $\ord z_k,\ord w_k>-2$ for all $k$, as nothing is circled
(Rule~1a at the outer vertices; on $T$ the strokes of the $zw$-edges are not maximal by
Estimates~1 and~2, so Rule~1e applies); and $\ord Z_{kl},\ord W_{kl}>-2$
unless $k,l\in T$.

Put $M_T=m_3+m_4+m_5$, $M=\sum_km_k$, $z_c=\sum_{k\in T}m_kz_k/M_T$, $w_c$ likewise,
$\xi_k=z_k-z_c$ and $\eta_k=w_k-w_c$ ($k\in T$), and let
$F_j=\sum_{k\in T}m_k\delta_{jk}^3(z_j-z_k)$ and $G_j=\sum_{k\in T}m_k\delta_{jk}^3(w_j-w_k)$ be the
forces of the triangle on the outer body $j$.  By Proposition~\ref{prop:d45-first}, (23) holds
and $(\xi/t,\eta/t)$ tends to a zero-multiplier triangle.  Renumbering $T$,
$(\sqrt{m_3},\sqrt{m_4},\sqrt{m_5})=\lambda(b,b(b+1),b+1)$ with $b,\lambda>0$.
This triangle has real, mass-centred, pairwise distinct offsets $s_3,s_4,s_5$ with distinct
cubes; $s_4,s_5\ne0$, and $s_3=0$ exactly for $b=1$, the \emph{central-body slice} $m_4=m_5=4m_3$.
Its $z$- and $w$-shapes agree, and it is unique up to translations, the scalings
$(\xi,\eta)\mapsto(a\xi,a'\eta)$ and the sign of its inverse distances.  A constant
reparametrization and a constant renormalization $(z,w)\mapsto(az,a^{-1}w)$, which preserves
rows, ties, orders, the $r_{kl}^2$, $U$ and $I$, give
\begin{equation}\label{d5-eq-offsets}
  \xi_k=t\,(s_k+\sigma_k),\qquad \eta_k=t\,(s_k+\tau_k),\qquad \sigma_k,\tau_k\to0.
\end{equation}

For $j\in\{1,2\}$ let $\alpha_{jc},\beta_{jc}>-2$ be the orders and $\hat z_j,\hat w_j$ the leading
coefficients of $z_j-z_c$, $w_j-w_c$; for $P\in\{12,1c,2c\}$ let $\alpha_P,\beta_P$ be the orders of
the $z$- and $w$-differences, $r_P^2$ their product and $\rho_P=\frac12(\alpha_P+\beta_P)=\ord r_P$.
The \emph{pattern} of $\gamma$ is the word $X_{12}X_{1c}X_{2c}$, where $X_P$ is C, R or E as
$\rho_P>0$, $=0$ or $<0$.

\subsubsection{The reality input}\label{d5-reality}
\begin{lemma}\label{d5-lem-reality}
Let $\Gamma\subset A_m$ be an irreducible curve containing a real point $q$.  Then $U=I=I_0$ on
$\Gamma$, where $I_0=\sum_km_k|z_k(q)|^2>0$, and along every Diagram 5 germ of $\Gamma$, in the
frame above,
\begin{equation}\label{d5-eq-inertia}
  I_T+I_3=I_0,\qquad I_T=\sum_{k\in T}m_k\xi_k\eta_k,\qquad
  I_3=\frac{m_1m_2r_{12}^2+m_1M_Tr_{1c}^2+m_2M_Tr_{2c}^2}{M},
\end{equation}
with $\ord I_T=2$ and hence $\ord I_3=0$.
\end{lemma}

\begin{proof}
By Lemma~\ref{lem:U-constant}, $U$ is constant on $\Gamma$, and $U=I$ on $A_m$.  At $q$,
$w_k=\bar z_k$ and $z_1(q)\ne z_2(q)$ (tie of $12$), so $I(q)>0$.  The renormalization preserves
$U$ and $I$.  As $\sum_km_kz_k=\sum_km_kw_k=0$ on $A_m$, $I-I_T-I_3$ vanishes identically, and
$I_T=t^2(\sum_km_ks_k^2+o(1))$ with $\sum_km_ks_k^2>0$ by \eqref{d5-eq-offsets}.
\end{proof}

This is the only input from reality.  It enters through the conditions (A) and (B) of
\S\ref{d5-census}, that two quantities have order exactly $0$; so only $I_0\ne0$ is used, not its
sign, nor $\sigma(\Gamma)=\Gamma$.

\subsubsection{Regimes of an outer body}\label{d5-regimes}
\begin{definition}\label{d5-def-regimes}
Let $j\in\{1,2\}$, $\alpha=\alpha_{jc}$ and $\beta=\beta_{jc}$.  The outer body $j$ is
\begin{alphalist}
\item \emph{far} if $\alpha,\beta<1$.  Then $\delta_{jk}/\delta_{jc}\to\varepsilon_{jk}=\pm1$ for a
fixed branch $\delta_{jc}$ of $((z_j-z_c)(w_j-w_c))^{-1/2}$.  The body is a \emph{monopole} if
$\mu_j=\sum_km_k\varepsilon_{jk}\ne0$.  Otherwise it is \emph{dipole-$z$} if $\alpha>\beta$,
\emph{dipole-$w$} if $\beta>\alpha$, and, if $\alpha=\beta$, \emph{dipole-equal-$F$},
\emph{dipole-equal-$G$} or \emph{dipole-equal} as $\hat w_j+3\hat z_j$, $3\hat w_j+\hat z_j$ or
neither vanishes, the leading coefficients of $F_j$, $G_j$ at the monopole orders raised by
$1-\alpha$ being nonzero multiples of these.
\item \emph{near in $z$} if $\alpha\ge1>\beta$.  For $\alpha=1$ let
$\theta_j=\lim t^{-1}(z_j-z_c)$.  The body is \emph{$z$-dominant-coincidence} if $\alpha=1$ and
$\theta_j=s_k$, its \emph{special body} $k$, so that $\ord(z_j-z_k)=1+e_j$ with $e_j>0$;
\emph{$z$-coincident} if $\alpha=1$ and $\theta_j\notin\{s_3,s_4,s_5\}$, with the variants
\emph{$z$-coincident-$\Phi$}, \emph{-$\Psi$}, \emph{-both} if $\Phi_j=0\ne\Psi_j$,
$\Psi_j=0\ne\Phi_j$, $\Phi_j=\Psi_j=0$, where $\Phi_j=\sum_km_ky_{jk}$, $\Psi_j=\sum_km_ky_{jk}^3$,
$\delta_{jk}=y_{jk}\nu_j(1+o(1))$, $y_{jk}^{-2}=\theta_j-s_k$ and $\nu_j$ is a branch of
$(t(w_j-w_c))^{-1/2}$; \emph{$z$-inner} if $\alpha>1$ off the slice; and
\emph{$z$-dominant-central} if $\alpha>1$ on the slice, with special body $3$,
$1+e_j=\ord(z_j-z_3)$ and $g=\ord\xi_3$.
\item \emph{near in $w$} if $\beta\ge1>\alpha$, with the seven mirror regimes.
\end{alphalist}
\end{definition}

\begin{lemma}\label{d5-lem-complete}
Let $m\in\Rp^5$.  Each outer body lies in exactly one of these $20$ regimes, and
$z$-coincident-both and $w$-coincident-both are empty.  The dipoles need $m_a=m_b+m_c$ for the
minority body $a$, which fails on the slice; and for a $z$-dominant-central body the least of
$\alpha_{jc}$, $1+e_j$, $g$ is attained at least twice.
\end{lemma}

\begin{proof}
For $k\in T$, $z_j-z_k=(z_j-z_c)-\xi_k$ and $w_j-w_k=(w_j-w_c)-\eta_k$, with
$\ord\xi_k,\ord\eta_k\ge1$.  If $\alpha,\beta\ge1$, then $\ord r_{jk}\ge1$ and $jk$ would be a
$zw$-edge.  If $\alpha,\beta<1$, then $\delta_{jk}^2/\delta_{jc}^2\to1$; $\mu_j=0$ needs a single
minority body $a$ with $m_a=m_b+m_c$, false for the slice masses $(1,4,4)$; and
$\hat w_j+3\hat z_j=3\hat w_j+\hat z_j=0$ would give $\hat z_j=\hat w_j=0$.  Let $\alpha=1>\beta$:
at most one $s_k$ equals $\theta_j\ne0$, and otherwise every $z_j-z_k$ has order $1$; at positive
masses on (23), $\Phi_j=\Psi_j=0$ forces $y=0$, i.e.\ $\theta_j=\infty$.  Let $\alpha>1>\beta$: then
$z_j-z_k=-ts_k(1+o(1))$ when $s_k\ne0$; on the slice $\ord(z_j-z_3)\ge\min(\alpha,g)>1$, and
$z_j-z_c=(z_j-z_3)+\xi_3$ gives the tropical condition, including the deep case
$\alpha=g<1+e_j$.  The rest is symmetric.
\end{proof}

If $z_j\equiv z_c$ or $\xi_3\equiv0$, the order $\alpha_{jc}$ or $g$ is infinite.  Such an order
enters the constraints below only through conditions that remain true when it is replaced by any
sufficiently large finite value, and the census, whose variables are free, includes such values.
The order
$\alpha_{12}$ is always finite, as the tie of $12$ gives $z_1\ne z_2$.  Up to the exchanges
$1\leftrightarrow2$ and $z\leftrightarrow w$, which permute patterns and four-products, there are
$111$ pairs of regimes.

\subsubsection{Identities and orders}\label{d5-census}
Let $\Pi=\delta_{12}^3(z_1-z_2)$, $U_T=\sum_{k<l\in T}m_km_l\delta_{kl}$ and, summing over
$j\in\{1,2\}$, $k\in T$, $V_1=\sum m_jm_k\delta_{jk}^3\xi_k(w_k-w_j)$,
$V_2=\sum m_jm_k\delta_{jk}^3\eta_k(z_k-z_j)$.  The census uses the conditions below; the identities
in (R), (T), the virial identity $I_T=U_T+\frac12(V_1+V_2)$ and the pair-centre identity of
Proposition~\ref{d5-prop-over} are exact combinations of rows and ties, as a symbolic computation
shows (Appendix~\ref{app:verification}).
\begin{itemize}
\item[(R)] $z_1-z_2=(m_1+m_2)\Pi+F_1-F_2$, $z_1-z_c=m_2\Pi+(1+\frac{m_1}{M_T})F_1+\frac{m_2}{M_T}F_2$
and $z_2-z_c=-m_1\Pi+\frac{m_1}{M_T}F_1+(1+\frac{m_2}{M_T})F_2$, with their $w$-analogues; the three
forms in $(\Pi,F_1,F_2)$ are pairwise independent on each coordinate plane, with common kernel
$(1,-m_2,m_1)$;
\item[(C)] $z_1-z_2=(z_1-z_c)-(z_2-z_c)$, and likewise in $w$;
\item[(T)] $V_1=V_2$: the outer bodies exert no torque on the triangle about its centre;
\item[(A)] $\ord I_3=0$ (Lemma~\ref{d5-lem-reality});
\item[(B)] $\sum_{j,k}m_jm_k\delta_{jk}\bigl(1+\frac12\bigl(\frac{\xi_k}{z_j-z_k}+\frac{\eta_k}{w_j-w_k}\bigr)\bigr)
+m_1m_2\delta_{12}=I_0-I_T$, of order exactly $0$, from $U=I_0$ and the virial identity;
\item[(V)] $V_1=V_2=I_T-U_T$, where $\ord I_T=2$ and $\ord U_T\ge\min(5,-1+2\kappa)$ by
\eqref{d5-eq-UT};
\item[(M)] for two outer bodies near in the same colour, the relations of
Lemma~\ref{d5-lem-consequences}.
\end{itemize}

Each regime fixes the orders of $F_j$, $G_j$, of the shares of body $j$ in the torque
$V_1-V_2$, in (B), $V_1$ and $V_2$, of its forcings of the triangle bodies and of the $r_{jk}^2$:
exactly with nonzero leading coefficient, by lower bounds, by the vanishing of the generic
coefficient (of $F_j$ in dipole-equal-$F$ and $z$-coincident-$\Phi$, of $G_j$ in dipole-equal-$G$
and $z$-coincident-$\Psi$), or not at all.  Nonvanishing comes from $\mu_j\ne0$, from the dipole
moment $\sum_km_k\varepsilon_{jk}\xi_k=-2m_a\xi_a$ ($a$ the body of minority sign; of order $1$ unless
$a$ is central), from $\sum_km_ks_k^2>0$, from the definitions, from $s_k\ne0$ at a dominant special
body, and, for $z$-inner, from the following observation: off the slice a single $s_k$ has the
minority sign, and as $y_{jk}^{-2}=-s_k$, its term alone forms the real or the imaginary part of
$\Phi_j$ and of $\Psi_j$.  A second computation re-derives these orders in the $12$ regimes of the $z$-side ($w$ by
symmetry) from exact first-order expansions with remainder bounds, and checks them, every census
constraint and the formulas for distances and four-products by exact linear programming over each
regime; it finds one exception (\S\ref{d5-slice}).  Valid constraints that the census omits, such
as $\alpha_{jc},\beta_{jc}>-2$, only weaken it.

A vanishing sum of Puiseux series attains its least order twice; a sum of order exactly $0$ has
least order $\le0$, attained twice if negative; and in (R), at the least order of $\Pi,F_1,F_2$, at
most one equation cancels if two of them attain it, and none, one or all three if all do.  So
each identity becomes a finite disjunction of rational linear systems, with strict and non-strict
inequalities, in the orders.  The census explores the disjuncts depth first and decides each
system exactly over $\mathbb Q$ by the simplex method, with free variables and a common slack
$0<s\le1$ for the strict inequalities.

\subsubsection{The Morse--Bott lemma}\label{d5-mb}
In reduced coordinates $\hat x_k=\xi_k/t$, $\hat y_k=\eta_k/t$, $D_{kl}=t\delta_{kl}$ ($k,l\in T$)
we have $U_T=t^{-1}\hat U$, $\hat U=\sum_{k<l}m_km_lD_{kl}$, with
$D_{kl}=((\hat x_k-\hat x_l)(\hat y_k-\hat y_l))^{-1/2}$ on the branch of the limit.  For
$\psi_k=\sum_lm_lD_{kl}^3(\hat y_k-\hat y_l)$ and $\phi_k=\sum_lm_lD_{kl}^3(\hat x_k-\hat x_l)$,
$\partial\hat U/\partial\hat x_k=-\frac{m_k}2\psi_k$, $\partial\hat U/\partial\hat y_k=-\frac{m_k}2\phi_k$
and, by homogeneity, $\hat U=\sum_km_k\hat x_k\psi_k=\sum_km_k\hat y_k\phi_k$.  The triangle rows
read exactly
\begin{equation}\label{d5-eq-rows}
  \psi_k=t^3\hat y_k-\hat g_k,\qquad \phi_k=t^3\hat x_k-\hat g^z_k,
\end{equation}
where $g_k=t^2\sum_jm_j\delta_{jk}^3(w_k-w_j)$ is the crossed forcing of the $w$-row of $k$,
$g^z_k$ its $z$-analogue, and the hat is the deviation from the $m$-weighted mean over $T$.  Let
$O=\{(as+u\mathbf 1,a's+v\mathbf 1)\}$ be the orbit of the zero-multiplier triangle under
translations and scalings; the reduced triangle tends to $x_0=(s,s)\in O$.

\begin{lemma}[computer-assisted]\label{d5-lem-mb}
At each point $(as,a's)$ of $O$, $\hat U=0$, $d\hat U=0$, the translations and scalings lie in
the Hessian kernel, and on $(\hat x_3,\hat y_3)$ the Hessian $H$ has entries $3X/(\varrho a^2)$,
$X/(\varrho aa')$, $3X/(\varrho a'^2)$, where $\varrho^2=aa'$ and $X=b^4(b+1)^6/4$ for masses
normalized to $b^2,b^2(b+1)^2,(b+1)^2$; its determinant is $8X^2/\varrho^6$, and
$(H^{-1})_{\hat x\hat x}=3a^2\varrho/(8X)\ne0$.
\end{lemma}
The proof is an exact computation in the field $\mathbb Q(b,a,\varrho)$ of rational functions
(Appendix~\ref{app:verification}).

\begin{lemma}[Morse--Bott]\label{d5-lem-mb2}
If $x(t)\to x_0$ is a germ of reduced triangles, $\kappa_\psi=\ord\psi(x(t))$ and
$\kappa_\phi=\ord\phi(x(t))$, then $\ord\hat U(x(t))\ge2\min(\kappa_\psi,\kappa_\phi)$, and
$\kappa_\psi<\kappa_\phi$ implies $\kappa_\phi\le2\kappa_\psi$ (and symmetrically).
\end{lemma}

\begin{proof}
As $s_4\ne s_5$, the $\hat x_3$- and $\hat y_3$-directions span a complement of the tangent space of
$O$, which by Lemma~\ref{d5-lem-mb} is thus the Hessian kernel; write points near $x_0$ as
$o(u)+v_1e^x_3+v_2e^y_3$, $o(u)\in O$.  Taylor's formula gives $\hat U=\frac12v^TH(u)v+O(|v|^3)$ and
$\partial_v\hat U=H(u)v+O(|v|^2)=-\frac{m_3}2(\psi_3,\phi_3)$ with $H(u)$ invertible, so
$\ord\hat U\ge2\ord v=2\ord\partial_v\hat U\ge2\min(\kappa_\psi,\kappa_\phi)$.  If
$\kappa_\psi<\kappa_\phi$, the leading coefficient of $\psi$ satisfies $\sum_km_k\psi_k=0$ and, as
$\hat U=\sum_km_k\hat x_k\psi_k$ has order $>\kappa_\psi$, $\sum_km_ks_k\psi_k=0$; as $s_4\ne s_5$,
its third component is nonzero, so $\partial_v\hat U=t^{\kappa_\psi}(\chi,0)+o(t^{\kappa_\psi})$ with
$\chi\ne0$ and $\hat U=\frac12(\partial_v\hat U)^TH^{-1}\partial_v\hat U+O(|v|^3)
=\frac12(H^{-1})_{\hat x\hat x}\chi^2t^{2\kappa_\psi}+o(t^{2\kappa_\psi})$.  But
$\hat U=\sum_km_k\hat y_k\phi_k$ has order $\ge\kappa_\phi$.
\end{proof}

By Lemma~\ref{d5-lem-mb2}, as the self-terms in \eqref{d5-eq-rows} have order $3$, along every
Diagram~5 germ
\begin{equation}\label{d5-eq-UT}
  \ord U_T\ge\min(5,-1+2\kappa),\qquad \kappa=\min(\ord\hat g,\ord\hat g^z).
\end{equation}
\begin{corollary}[uniform forcing]\label{d5-cor-uniform}
Let both outer bodies be near in $z$, and let $\kappa_*$ be the least order of the terms
$t^2m_j\delta_{jk}^3(w_k-w_j)$ ($j\in\{1,2\}$, $k\in T$).  Then the coefficient of $t^{\kappa_*}$ in
$g_k$ is the same for $k=3,4,5$; likewise with $z$ and $w$ exchanged.
\end{corollary}

\begin{proof}
Let $\ord(z_j-z_k)=1+e_{jk}$, $e_{jk}\ge0$.  The term of body $j$ in $g_k$ has order
$\kappa_{jk}=\frac12(1-\beta_{jc}-3e_{jk})$, in $g^z_k$ the order
$\frac12(3-3\beta_{jc}-e_{jk})=2\kappa_{jk}+\frac12(1-\beta_{jc}+5e_{jk})$.  As $\beta_{jc}>-2$ and
$e_{jk}=0$ for some $k$, $\kappa_*<\frac32$.  Unequal coefficients would give $\kappa_\psi=\kappa_*$
(as $3>\kappa_*$), while $\hat g^z$ has order $>2\kappa_*$ (as $\beta_{jc}<1$) and $3>2\kappa_*$; so
$\kappa_\phi>2\kappa_\psi$, with $\kappa_\psi>0$ as $\psi\to0$, against Lemma~\ref{d5-lem-mb2}.
\end{proof}

\begin{lemma}[computer-assisted]\label{d5-lem-consequences}
Let both outer bodies be near in $z$ (or, mirrored, in $w$).  (a) A forcing from one source,
$k\mapsto c(\theta-s_k)^{-3/2}$ with $c\ne0$, or supported on fewer than three triangle bodies, is
not uniform.  (b) Two $z$-dominant-coincidence bodies share their special body; so
$\alpha_{12}>1$.  (c) A $z$-coincident-$\Phi$ or $z$-coincident-$\Psi$ body does not occur with a
$z$-dominant-coincidence body.  (d) A $z$-inner body does not occur with a $z$-dominant-central
body.  (e) Two $z$-dominant-central bodies have the same $g$; their forcings of body $3$, of orders
$\varphi_j=\ord(\delta_{j3}^3(w_3-w_j))=-\frac32(1+e_j)-\frac12\beta_{jc}$, tie and cancel;
$g>3+\varphi_1$; and $1+e_j=g<\alpha_{jc}$ with $\beta_{jc}<-g$ for both $j$ is impossible.
\end{lemma}
The lemma follows from Corollary~\ref{d5-cor-uniform} by exact computations:
(a) through the resultant $9(s_4-s_5)^2\cdot\frac12\sum_{k<l}(s_k-s_l)^2\ne0$ of
$(\theta-s_3)^3-(\theta-s_4)^3$ and $(\theta-s_3)^3-(\theta-s_5)^3$, (c) through twelve gcd
computations over $\mathbb Q(\omega)$ at the points $(s_k-\omega s_{k'})/(1-\omega)$,
$\omega^3=1\ne\omega$, followed by exact real-root counting in $b$ (Appendix~\ref{app:verification});
(d) is the definition, and the last case of (e) falls as in Proposition~\ref{d5-prop-over}.

\begin{proof}[Proof of the tie and of $g>3+\varphi_1$ in (e)]
The term of body $j$ in $g_k$ has order $2+\varphi_j$ for $k=3$, and order
$\frac12-\frac12\beta_{jc}>2+\varphi_j$ for $k=4,5$, where $\ord(z_j-z_k)=1$ and $e_j>0$.  So the
leading crossed forcing is supported on body $3$, and by Corollary~\ref{d5-cor-uniform} and (a)
the two forcings of body $3$ have one order $\varphi:=\varphi_1=\varphi_2$ and cancel; hence
$\ord\hat g_3>2+\varphi$.  The $z$-forcing of body $3$ by body $j$ has order
$2+\varphi+(1+e_j)-\beta_{jc}>2+\varphi$, as $\beta_{jc}<1$, and the self-terms in
\eqref{d5-eq-rows} have order $\ge3>2+\varphi$, as $\varphi<-\frac12$ by $\beta_{jc}>-2$.  So
$\ord\psi_3,\ord\phi_3>2+\varphi$.  In the proof of Lemma~\ref{d5-lem-mb2},
$\partial_v\hat U=H(u)v+O(|v|^2)=-\frac{m_3}2(\psi_3,\phi_3)$ with $H(u)$ invertible and $v\to0$,
so $\ord v\ge\min(\ord\psi_3,\ord\phi_3)>2+\varphi$.  On the slice $s_3=0$, so the mass-centred
point $o(u)$ of the orbit has third $\hat x$-component $0$, and mass-centring gives
$\hat x_3=(1-m_3/M_T)v_1$.  Therefore $g=\ord\xi_3=1+\ord\hat x_3\ge1+\ord v>3+\varphi$.
\end{proof}

\subsubsection{The census and over-rotation}\label{d5-result}

\begin{lemma}[computer-assisted]\label{d5-lem-census}
For each of the $111$ pairs of regimes and each of the $27$ patterns, the conditions of the frame
and of the two regimes together with (R), (C), (T), (A), (B), (V), (M) are infeasible, except for
the pairs and patterns of Table~\ref{d5-tab-survivors} and for pairs containing an empty regime.
For the five monopole pairs of \S\ref{d5-slice} this holds with the rows corrected there.
\end{lemma}
The proof is a finite computation: for each of the $111$ pairs of regimes and each of the $27$
patterns, the finite disjunction of rational linear systems obtained above is decided by exact
linear programming over $\mathbb Q$, with the orders checked by the second computation and, for
the five monopole pairs, with the corrected rows (Appendix~\ref{app:verification}).

\begin{table}[ht]
\centering
\small
\begin{tabular}{p{8.9cm}p{3.3cm}}
\hline
pair of regimes (up to the exchanges) & surviving patterns \\
\hline
(monopole, monopole); (monopole, $X$) for $X$ = dipole-$z$, dipole-equal, dipole-equal-$F$;
(dipole-$z$, dipole-$z$); (dipole-equal, $X$) for $X$ = dipole-equal, dipole-equal-$F$;
(dipole-equal-$F$, $X$) for $X$ = dipole-equal-$F$, dipole-equal-$G$ & RRR \\
($X$, $Y$) for $X,Y\in\{z$-coincident, $z$-coincident-$\Phi$, $z$-coincident-$\Psi\}$ & RRR \\
($z$-inner, $z$-inner) & RRR \\
($z$-inner, $X$) for $X\in\{z$-coincident, $z$-coincident-$\Phi$, $z$-coincident-$\Psi\}$ & RCR \\
($z$-dominant-coincidence, $z$-dominant-coincidence) & CEE, CRR \\
($z$-dominant-central, $z$-dominant-central) & RRR, RRC, RCR, CEE, CRR \\
every other pair of nonempty regimes & none \\
\hline
\end{tabular}
\caption{The census of Lemma~\ref{d5-lem-census}.  A pattern lists the signs of $\rho_{12}$,
$\rho_{1c}$, $\rho_{2c}$: C positive, R zero, E negative.}\label{d5-tab-survivors}
\end{table}

Each four-product contains a triangle edge, of order $2$, and its order is the sum of its edge
orders: $\alpha_{12}+\beta_{12}$ for $r_{12}^2$; for $r_{jk}^2$, $\alpha_{jc}+\beta_{jc}$ if $j$ is far
and $1+\beta_{jc}$ (or $1+e_j+\beta_{jc}$ at the special body) if $j$ is near in $z$.

\begin{lemma}[computer-assisted]\label{d5-lem-four}
For every pair and pattern of Table~\ref{d5-tab-survivors} other than ($z$-inner, $z$-inner) with
RRR, adding $\ord p_{ab}<0$ for any one of the ten four-products makes the constraints infeasible.
For ($z$-inner, $z$-inner) with RRR exactly the six four-products $p_{jk}$, $j\in\{1,2\}$, $k\in T$,
can have negative order.
\end{lemma}
The proof is a computation of the same kind: for each surviving pair and pattern and each of the
ten four-products, $260$ cases in all, we add $\ord p_{ab}<0$, with $\ord p_{ab}$ expressed through
the edge orders as above, and decide the resulting systems by exact linear programming over
$\mathbb Q$ (Appendix~\ref{app:verification}).

By Lemmas~\ref{d5-lem-census} and~\ref{d5-lem-four} (and \S\ref{d5-slice} for pairs with a
monopole), a germ with an unbounded four-product would be \emph{over-rotated}: both outer bodies
inner in one colour, say $z$, and all $\rho_P=0$, so that $\beta_{jc}=-\alpha_{jc}<-1$ and the
crossed distances, of order $\frac12(1+\beta_{jc})<0$, are unbounded.

\begin{proposition}\label{d5-prop-over}
Let $m\in\Rp^5$.  No Diagram 5 germ has both outer bodies $z$-inner and pattern RRR, nor both
$w$-inner and pattern RRR.
\end{proposition}

\begin{proof}
Here $\beta_{jc}=-\alpha_{jc}<-1$ and every $z_j-z_k$ ($k\in T$) has order $1$; as
$\delta_{jk}^3(w_j-w_k)=\delta_{jk}/(z_j-z_k)$,
$\ord G_j\ge-\frac32-\frac12\beta_{jc}=\beta_{jc}+\frac32(\alpha_{jc}-1)>\beta_{jc}$.  By (R) in $w$, the
pair-centre identity $m_1(w_1-w_c)+m_2(w_2-w_c)=\frac M{M_T}(m_1G_1+m_2G_2)$ holds.  If
$\beta_{1c}<\beta_{2c}$ (the other case is symmetric), its left side has order $-\alpha_{1c}$ and its
right side order at least
$-\frac32+\frac12\alpha_{2c}$, larger by $(\alpha_{1c}-1)+\frac12(\alpha_{2c}-1)>0$.  So
$\beta_{1c}=\beta_{2c}=:\beta$ and $m_1\hat w_1+m_2\hat w_2=0$.  Next, $z_k-z_j=ts_k(1+o(1))$ and
$w_k-w_j=-\hat w_jt^\beta(1+o(1))$, so
$(t^2m_j\delta_{jk}^3(w_k-w_j))^2=-(m_j^2/\hat w_j)s_k^{-3}t^{1-\beta}(1+o(1))$, and the coefficient of
$t^{(1-\beta)/2}$ in $g_k$ is $(c_1\varepsilon'_{1k}+c_2\varepsilon'_{2k})s_k^{-3/2}$, with fixed
branches, $c_j^2=-m_j^2/\hat w_j$ and signs $\varepsilon'_{jk}$.  By Corollary~\ref{d5-cor-uniform}
it is one value $f$ for $k=3,4,5$.  If $f\ne0$, then $f^2s_k^3=(c_1\varepsilon'_{1k}+c_2\varepsilon'_{2k})^2$
depends only on $\varepsilon'_{1k}\varepsilon'_{2k}$, so two of the cubes $s_k^3$ coincide, which is
false.  Hence $f=0$ and $c_1^2=c_2^2$, that is, $m_1^2\hat w_2=m_2^2\hat w_1$.  The two linear
relations for $(\hat w_1,\hat w_2)$ have determinant $m_1^3+m_2^3>0$, so $\hat w_1=\hat w_2=0$,
which is absurd.
\end{proof}
No reality input is used here.

\subsubsection{The central-body slice}\label{d5-slice}
On the slice body $3$ sits at the centre of the zero-multiplier triangle; this gives the
dominant-central regimes, voids $z$-inner, and causes the one exception found by the second
computation: for a monopole whose minority crossed sign is on body $3$ the dipole moment
$-2m_3\xi_3$ has order $g>1$, so for its shares in $V_1$, $V_2$ the orders the census uses as exact
are only lower bounds (the coefficients there vanish).  Its other rows are those of every
monopole: the leading coefficients of $F_j$, $G_j$ and of its term in (B) are multiples of
$\mu_j\ne0$, and its torque row is a lower bound.  The correction can matter only for a monopole
facing a monopole, a $z$-coincident (or -$\Phi$, -$\Psi$) or a $z$-dominant-coincidence body
(elsewhere the masses avoid the slice, the regime is void on it, or there is no virial row);
recomputed with it, these five pairs keep their surviving patterns.

\begin{remark}\label{d5-lem-ablation}
Suppose that we drop (T), (B), (A), (V), and keep of (M) only its exact part ($z$-inner is void
on the slice; two $z$-dominant-central bodies share $g$ and satisfy the tropical condition of
$z_1-z_2=(z_1-z_3)-(z_2-z_3)$).  Then, among all pairs of nonempty regimes, patterns and
four-products, a four-product of negative order is feasible only for ($z$-inner, $z$-inner) with
RRR, ($z$-dominant-coincidence, $z$-dominant-coincidence) with EEE or REE, and
($z$-dominant-central, $z$-dominant-central) with RRR, RRC, RCR, CEE, CRR or CCC.  This shows
where (A), (B) and the Morse--Bott relations are needed.  None of these pairs contains a monopole,
so the slice exception cannot create an unbounded four-product even before the correction.  We do
not use this.
\end{remark}

\subsubsection{Proof of Theorem~\ref{thm:d5}}
Let $\gamma$ be a Diagram 5 germ with an unbounded four-product on an irreducible curve
$\Gamma\subset A_m$ containing a real point.  By Lemma~\ref{d5-lem-reality}, (A) and (B) hold along
$\gamma$; the other conditions of \S\ref{d5-census} hold along every germ, and all of them are
necessary conditions on the orders.  By Lemma~\ref{d5-lem-complete} and the exchanges, the outer
bodies form one of the $111$ pairs of regimes, both nonempty, and some $p_{ab}$ has negative order
(Definition~\ref{def:germ}).  If the pair contains a monopole, Lemma~\ref{d5-lem-census}, with the
rows of \S\ref{d5-slice} for the five slice-sensitive pairs, leaves only pairs of two far bodies
with pattern RRR (Table~\ref{d5-tab-survivors}).  Then $\rho_{12}=\rho_{1c}=\rho_{2c}=0$, so
$r_{12}^2$ and every $r_{jk}^2$ with $j\in\{1,2\}$, $k\in T$ have order $0$, and every four-product,
which contains a triangle edge of order $2$, has order at least $2$: a contradiction.  Otherwise
every order used by the census is confirmed by the second computation, and
Lemmas~\ref{d5-lem-census} and~\ref{d5-lem-four} leave only ($z$-inner, $z$-inner) with RRR or
its mirror, which Proposition~\ref{d5-prop-over} excludes.  \qed

\begin{remark}\label{d5-rem-h1}
We do not prove the \emph{outer-triple regularity} of Diagram~5 germs, that is, that their pattern
is RRR: the census leaves RCR, CEE, CRR and, on the slice, RRC, RCR, CEE, CRR.
Theorem~\ref{thm:d5} does not use it.
\end{remark}
